% SEGMENT OLP-0668-S01
% Part: proof-theory
% Chapter: natural-deduction
% Section: rules-proofs

\documentclass[../../../include/open-logic-section]{subfiles}

\begin{document}

\olfileid{pt}{nat}{rp}

\olsection{விதிகளும் \usetoken{P}{proof}களும்}

சில வரையறைகளைக் கொடுத்து, பயன்படுத்தவுள்ள சொற்களின்
பொருளைத் தெளிவுபடுத்துவதிலிருந்து தொடங்குவோம்.

\Log{N1c}, \Log{N1i} ஆகியவற்றில் $!A \lif !B$ ஐக் கொண்டிருக்கும்
பின்வரும் இரு விதிகளைக் கவனிப்போம்.
\[
\bottomAlignProof
\AxiomC{$\Discharge{!A}{x}$}
\shortDeduce
\DeduceC{$!B$}
\DischargeRule{\Intro{\lif}}{x}
\UnaryInfC{$!A \lif !B$}
\DisplayProof
\qquad
\bottomAlignProof
\AxiomC{$!A \lif !B$}
\AxiomC{$!A$}
\RightLabel{\Elim{\lif}}
\BinaryInfC{$!B$}
\DisplayProof
\]
\Elim{\lif} விதி எளிமையானது. கோட்டுக்கு மேலுள்ள !!{formula}s
\emph{முன்கோள்கள்}. இடப்புற !!{formula} என்பது \emph{முதன்மை}
!!{formula}~$!A \lif !B$. ஒவ்வொரு !!{elimination} விதிக்கும்
இவ்வாறான ஒரு முன்கோள் உண்டு; அது எப்போதும் மிக இடப்புறத்தில்
இருக்கும். அதனை அந்த அனுமானத்தின் \emph{முதன்முன்கோள்} என்போம்.
இங்கு இருப்பதைப் போல் பிற முன்கோள்கள் இருந்தால், அவை
\emph{சிறுமுன்கோள்கள்}. !!{introduction} விதிகளின் முன்கோள்களும்
முதன்முன்கோள்கள் எனப்படும். கோட்டுக்குக் கீழுள்ள !!{formula}
\emph{முடிவு}.

% SEGMENT OLP-0668-S02
\Intro{\lif} விதிக்கு ஒரே ஒரு முன்கோள். இங்கு முடிவு முதன்மை
வாய்ப்பாடு~$!A \lif !B$. எல்லா !!{introduction} விதிகளிலும்
இவ்வாறே முடிவே முதன்மை வாய்ப்பாடு; அதன் முதன்மைச் செயலியே
``அறிமுகப்படுத்தப்படும்'' செயலியாகும்.

% SEGMENT OLP-0668-S03
$\Discharge{!A}{x}$ என்ற குறியீடு பின்வருவதைச் சொல்கிறது. ஓர்
!!a{proof} இல் \Intro{\lif} அனுமானம், முன்கோளான !!a{formula}~$!B$
மீது பயன்படுத்தப்படுகிறது. $!B$ இல் முடியும் உள்-!!{proof} இன்
தொடக்கத்தில் பல !!{formula}s உள்ளன. அவை \emph{எடுகோள்கள்}.
$!A$ வடிவிலான ஒன்று அல்லது பல எடுகோள்களிலிருந்து $!B$ இன்
!!^a{proof} என்பது $!A$ இலிருந்து $!B$ பின்தொடர்வதற்கான
!!a{proof}. \Intro{\lif} விதியைப் பயன்படுத்தும்போது $!A \lif !B$
என முடிக்கிறோம்; அப்போது முடிவு $!A$ என்ற எடுகோளை இனிச்
சார்ந்திருக்காது. இதைக் குறிக்க, \Intro{\lif} அனுமானத்தையும்
உள்ளடக்கிய !!{proof} இல் $!A$ வடிவிலான எடுகோள்கள்
\emph{விடுவிக்கப்படுகின்றன}, அல்லது \emph{மூடப்படுகின்றன}.
எந்த எடுகோள்களை விடுவிக்கலாம் என்பதை விதி குறிப்பிடுகிறது.
முடிவு $!A \lif !B$ ஆகும் \Intro{\lif} இல் அவை நிபந்தனைக்
கூற்றின் முற்பகுதியான $!A$ உடன் பொருந்தும் எடுகோள்கள்.
எந்த விதியில் எந்த எடுகோள்கள் விடுவிக்கப்படுகின்றன என்பதைக்
காட்டுவது சில நேரங்களில் முக்கியம்; அதற்காக விடுவிப்பு
அடையாளம்~$x$ பயன்படுகிறது. முக்கியமாக, எடுகோள்களை
விடுவிக்கும் விதிகள் தேவையான வடிவிலான \emph{எல்லா}
எடுகோள்களையும், அல்லது ஒன்றையும்கூட, விடுவிக்க
\emph{வேண்டியதில்லை}. எடுத்துக்காட்டாக, பின்வரும் நிறுவல்
சரியானதே; இதில் முதல் \Intro{\lif} எந்த எடுகோளையும்
விடுவிப்பதில்லை:
\[
\AxiomC{$\Discharge{!A}{y}$}
\RightLabel{\Intro{\lif}}
\UnaryInfC{$!B \lif !A$}
\DischargeRule{\Intro{\lif}}{y}
\UnaryInfC{$!A \lif (!B \lif !A)$}
\DisplayProof
\]
ஒரு விதி எந்த எடுகோளையும் விடுவிக்காதபோது, இங்கு
செய்ததுபோல் விடுவிப்பு அடையாளத்தை விடுகிறோம்.

% SEGMENT OLP-0668-S04
பின்வரும் !!{proof} இல்
\[
\AxiomC{$\Discharge{!A}{x}$}
\AxiomC{$\Discharge{!A}{y}$}
\RightLabel{\Intro{\land}}
\BinaryInfC{$!A \land !A$}
\RightLabel{\Elim{\land}}
\UnaryInfC{$!A$}
\DischargeRule{\Intro{\lif}}{x}
\UnaryInfC{$!A \lif !A$}
\DischargeRule{\Intro{\lif}}{y}
\UnaryInfC{$!A \lif (!A \lif !A)$}
\DisplayProof
\]
முதல் \Intro{\lif} அனுமானம் இடப்புற எடுகோள்~$!A^x$ ஐயும்,
இரண்டாவது வலப்புற எடுகோள்~$!A^y$ ஐயும் விடுவிக்கின்றன. அதாவது,
$x$ அடையாளமுள்ள எல்லா எடுகோள்களும் $x$ அடையாளமுள்ள
\Intro{\lif} ஆலும், $y$ அடையாளமுள்ளவை $y$ அடையாளமுள்ள
அனுமானத்தாலும் விடுவிக்கப்படுகின்றன. இது செயல்பட,
ஒரே அடையாளமுள்ள எல்லா எடுகோள்களும் ஒரே !!{formula}வாக
இருப்பது முக்கியம்.

% SEGMENT OLP-0668-S05
அளவையடை விதிகள் சற்றுச் சிக்கலானவை. எடுத்துக்காட்டாக,
\Log{N1i} இல் $\lexists$ க்கான விதிகள் பின்வருமாறு:
\[
\bottomAlignProof
\AxiomC{$!A(t)$}
\RightLabel{\Intro{\lexists}}
\UnaryInfC{$\lexists[x][!A(x)]$}
\DisplayProof
\qquad
\bottomAlignProof
\AxiomC{$\lexists[x][!A(x)]$} 
\AxiomC{$\Discharge{!A(c)}{x}$}
\shortDeduce
\DeduceC{$!B$}
\DischargeRule{\Elim{\lexists}}{x}
\BinaryInfC{$!B$} \DisplayProof 
\bottomAlignProof
\]
\Intro{\lexists} விதியில் முன்கோள்~$!A(t)$; இங்கு $t$ ஒரு மூடிய
சொல், அதாவது மாறிகள் இல்லாத சொல். முடிவு $\lexists[x][!A]$,
முன்கோள்~$\Subst{!A}{t}{x}$ ஆக அமையுமாறு ஏதோ ஓர்
!!{formula}~$!A$, மாறி~$x$, மூடிய சொல்~$t$ இருக்கும்போதும்
அப்போது மட்டுமே அந்த அனுமானம் சரியானது; அதாவது அது
\Intro{\lexists} திட்டவடிவ விதியுடன் பொருந்துகிறது.
\Elim{\lexists} விதி $!A(c)$ வடிவிலான எடுகோள்களை
விடுவிக்க அனுமதிக்கிறது: ஒவ்வோர் அனுமானத்திற்கும் ஏதோ ஒரு
!!{constant}~$c$ உண்டு; $\Subst{!A}{c}{x}$ வடிவிலான
எடுகோள்களை விடுவிக்கலாம். ஓர் அனுமானத்தில் விடுவிக்கப்படும்
எல்லா எடுகோள்களும் அதே~$c$ ஐப் பயன்படுத்த வேண்டும்.
இந்தக் குறியீட்டை \Elim{\lexists} அனுமானத்தின் \emph{தனிமாறி}
என்போம். விடுவிக்கப்படும்~$!A(c)$ எடுகோள்களைத் தவிர வேறு
எந்தத் திறந்த எடுகோளிலும் $c$ இடம்பெறாமலும், $!B$ இலும்
$!A$ இலும் இடம்பெறாமலும் இருந்தால் மட்டுமே அந்த அனுமானம்
சரியானது. இதுவே \emph{தனிமாறி நிபந்தனை}.

\Log{N1c}, \Log{N1i} ஆகியவற்றில் !!^{proof}s என்பவை
தொடக்க எடுகோள்களிலிருந்து அனுமான விதிகளால் தொகுத்தறிதலாகக்
கட்டப்படும் வாய்ப்பாடுகளின் முடிவுறு மரங்கள். ஒவ்வோர் இலையும்
ஓர் எடுகோள்; அதற்கு விடுவிப்பு அடையாளம்
இருக்கலாம். மற்ற ஒவ்வொரு !!{formula}வும் அதற்கு உடனே மேலுள்ள
!!{formula}s இலிருந்து முறைமையின் ஓர் அனுமான விதியால்
பின்வருகிறது. ஓர் அனுமானத்துக்கு மேலுள்ள மரத்தில் எடுகோளாக
வரும் !!{formula}வின் நிகழ்வு, அதற்கு மேலுள்ள அனுமானத்தால்
விடுவிக்கப்படாமல் இருந்தால், அந்த அனுமான நிலையில் அது
\emph{திறந்த} எடுகோள் எனப்படும்.

% SEGMENT OLP-0668-S06
\begin{defn}\ollabel{defn:proof}
\Log{N1c}, \Log{N1i} ஆகியவற்றின் \emph{!!^{proof}s} என்பவைப்
பின்வருமாறு தொகுத்தறிதலாக வரையறுக்கப்படும் முடிவுறு
!!{formula}s மரங்கள்:
\begin{enumerate}
  \item\ollabel{defn:prf:assumption} விடுவிப்பு அடையாளம்~$x$
  அல்லது $y$ இடப்பட்டிருந்தாலும் இல்லாவிட்டாலும், தனியாக
  நிற்கும் ஒவ்வொரு !!{formula}வும் !!a{proof}. ஒரே
  !!{formula}வை மட்டும் கொண்ட !!a{proof} இல் அந்த
  !!{formula} திறந்த எடுகோளாகக் கருதப்படும்.
  \item\ollabel{defn:prf:one-prem} $R$ என்பது ஒன்று, இரண்டு,
  அல்லது மூன்று முன்கோள்கள் $!A_1$, $!A_2$, $!A_3$ மற்றும்
  முடிவு~$!A$ உடைய விதியாகவும்; $\delta_1$, $\delta_2$,
  $\delta_3$ முறையே $!A_1$, $!A_2$, $!A_3$ இல் முடியும்
  !!{proof}s ஆகவும் இருந்தால், பின்வரும் வடிவங்கள்
  \[
  \AxiomC{}
  \RightLabel{$\delta_1$}
  \DeduceC{$!A_1$}
  \RightLabel{$R$}
  \UnaryInfC{$!A$}
  \DisplayProof
\qquad
  \AxiomC{}
  \RightLabel{$\delta_1$}
  \DeduceC{$!A_1$}
  \AxiomC{}
  \RightLabel{$\delta_2$}
  \DeduceC{$!A_2$}
  \RightLabel{$R$}
  \BinaryInfC{$!A$}
  \DisplayProof
\qquad
  \AxiomC{}
  \RightLabel{$\delta_1$}
  \DeduceC{$!A_1$}
  \AxiomC{}
  \RightLabel{$\delta_2$}
  \DeduceC{$!A_2$}
  \AxiomC{}
  \RightLabel{$\delta_3$}
  \DeduceC{$!A_3$}
  \RightLabel{$R$}
  \TrinaryInfC{$!A$}
  \DisplayProof
  \]
  முறையே !!a{proof} ஆகும்; அவற்றின் எடுகோள் அடையாளங்கள்
  பொருந்தியிருக்க வேண்டும் (வேறுபட்ட இரு !!{formula}s க்கு
  ஒரே விடுவிப்பு அடையாளம் இருக்கக்கூடாது), மேலும் அனுமானங்கள்
  $R$ இன் சரியான பயன்பாடுகளாக இருக்க வேண்டும்.
  \item புதிய !!{proof} இன் மிக மேலுள்ள !!{formula}
  நிகழ்வுகள் $\delta_1$, $\delta_2$, அல்லது~$\delta_3$ இன்
  திறந்த எடுகோள்களாக இருந்தால், அவை புதிய !!{proof} இன்
  \emph{திறந்த எடுகோள்கள்}; ஆனால் விதியால் விடுவிக்கப்படும்
  நிகழ்வுகள் விலக்கப்படும். விதிக்கு~$x$ (அல்லது $x\, y$)
  அடையாளம் இருந்தால், \Intro{\lif}, \Intro{\lnot} இல்
  $\delta_1$ இன் $x$ அடையாளமுள்ள எல்லாத் திறந்த எடுகோள்களும்;
  \Elim{\lexists} இல் $\delta_2$ இன் $x$ அடையாளமுள்ளவை
  அனைத்தும்; \Elim{\lor} இல் $\delta_2$ இன் $x$
  அடையாளமுள்ளவை அனைத்தும், $\delta_3$ இன் $y$
  அடையாளமுள்ளவை அனைத்தும் விடுவிக்கப்படும்.
\end{enumerate}
\end{defn}

% SEGMENT OLP-0668-S07
!!a{proof} இன் தொகுத்தறிதல் கட்டமைப்பில் பயன்படுத்தப்பட்ட
!!{proof}s அதன் \emph{உள்-!!{proof}s} எனப்படும். ஓர்
அனுமானத்தின் முன்கோள்களில் முடியும் உள்-!!{proof}s அந்த
அனுமானத்தின் \emph{உடனடி உள்-!!{proof}s} எனப்படும்.
\Log{N1c}, \Log{N1i} விதிகள் \olref{tab:N1} இல்
தரப்பட்டுள்ளன.

\subfile{rules-N1}

\begin{defn}
!!a{proof}~$\delta$ இன் \emph{!!{height}}~$\pheight{\delta}$
பின்வருமாறு தொகுத்தறிதலாக வரையறுக்கப்படுகிறது.
\begin{enumerate}
  \item $\delta$ ஓர் எடுகோளை மட்டும் கொண்டிருந்தால்,
  $\pheight{\delta} = 0$.
  \item $\delta$ ஒரு முன்கோளுடைய அனுமானத்தில் முடிந்து,
  அதன் உடனடி உள்-!!{proof}~$\delta_1$ என்றால்,
  $\pheight{\delta} = \pheight{\delta_1} + 1$.
  \item $\delta$ இரு முன்கோள்களுடைய அனுமானத்தில் முடிந்து,
  அதன் உடனடி உள்-!!{proof}s~$\delta_1$, $\delta_2$ என்றால்,
  $\pheight{\delta} = \max(\pheight{\delta_1}, \pheight{\delta_2}) + 1$.
  \item $\delta$ இன் இறுதி அனுமானம் \Elim{\lor} ஆகவும்,
  அதன் உடனடி உள்-!!{proof}s~$\delta_1$, $\delta_2$,
  $\delta_3$ ஆகவும் இருந்தால்,
  $\pheight{\delta} = \max(\pheight{\delta_1}, \pheight{\delta_2}, \pheight{\delta_3}) + 1$.
\end{enumerate}
தொடரணி~$S$ க்கு உயரம் $\le n$ உடைய !!a{proof}
இருந்தால், $\Proves[n] S$ என்று எழுதுவோம்.
\end{defn}

% SEGMENT OLP-0668-S08
\begin{ex}
\Log{N1i} இல் எந்த !!{formula}வும் தன்னையே திறந்த
எடுகோளாகக் கொண்டு !!a{proof} ஆகும். ஆகவே
\[
(!A \land !B)^x
\]
என்பது !!a{proof}~$\delta_1$. அதன் உயரம்~$0$.
$\delta_1$ இலிருந்து \Elim{\land} இன் இரு வடிவங்களில்
ஒவ்வொன்றைப் பயன்படுத்தி இரண்டு !!{proof}s~$\delta_2$,
$\delta_3$ ஐப் பெறலாம்:
\[
\AxiomC{$(!A \land !B)^x$}
\RightLabel{\Elim{\land}[2]}
\UnaryInfC{$!B$}
\DisplayProof
\qquad
\AxiomC{$(!A \land !B)^x$}
\RightLabel{\Elim{\land}[1]}
\UnaryInfC{$!A$}
\DisplayProof
\]
$\delta_2$, $\delta_3$ இன் இறுதி அனுமானங்களுக்கான
உடனடி உள்-!!{proof} இரண்டிலும் ஒன்றே: $!A \land !B$.
இரு !!{proof}s இலும் $!A\land !B$
நிகழ்வு திறந்த எடுகோளாக உள்ளது. உயரங்கள்
$\pheight{\delta_1} = 0$, $\pheight{\delta_2} =
\pheight{\delta_3} = 1$.

% SEGMENT OLP-0668-S09
\Intro{\land} விதி $!B$, $!A$ வடிவிலான இரு முன்கோள்களைக்
கொண்டு $!B \land !A$ வடிவிலான முடிவை நியாயப்படுத்துகிறது.
எனவே பின்வருவது !!a{proof}~$\delta_4$:
\[
\AxiomC{$(!A \land !B)^x$}
\RightLabel{\Elim{\land}[2]}
\UnaryInfC{$!B$}
\LeftSubproofLabel{$\delta_2$}
\AxiomC{$(!A \land !B)^x$}
\RightLabel{\Elim{\land}[1]}
\UnaryInfC{$!A$}
\RightSubproofLabel{$\delta_3$}
\RightLabel{\Intro{\land}}
\BinaryInfC{$!B \land !A$}
\DisplayProof
\]
அதன் உயரம் $\pheight{\delta_4} = \max(\pheight{\delta_2},
\pheight{\delta_3})+1 = \max(1,1) + 1 = 2$. அதற்கு
$!A \land !B$ இன் இரண்டு நிகழ்வுகள் எடுகோள்களாக உள்ளன;
இரண்டும் திறந்தவை.

இறுதியாக \Intro{\lif} ஐப் பயன்படுத்தி $\delta_5$ ஐப் பெறலாம்:
\[
\AxiomC{$\Discharge{!A \land !B}{x}$}
\RightLabel{\Elim{\land}[2]}
\UnaryInfC{$!B$}
\LeftSubproofLabel{$\delta_2$}
\AxiomC{$\Discharge{!A \land !B}{x}$}
\RightLabel{\Elim{\land}[1]}
\UnaryInfC{$!A$}
\RightSubproofLabel{$\delta_3$}
\RightLabel{\Intro{\land}}
\BinaryInfC{$!B \land !A$}
\DischargeRule{\Intro{\lif}}{x}
\UnaryInfC{$(!A \land !B) \lif (!B \land !A)$}
\DisplayProof
\]
அதன் உயரம் $\pheight{\delta_4} + 1 = 3$. \Intro{\lif}
அனுமானம் $x$ அடையாளமுடைய, $!A \land !B$ வடிவிலான
எடுகோள்களின் எல்லா நிகழ்வுகளையும், அதாவது உடனடி
உள்-!!{proof} இல் முன்பு திறந்திருந்த இரண்டு எடுகோள்களையும்
விடுவிக்கிறது. இவையே $\delta_5$ இன் ஒரே எடுகோள்கள்;
ஆகவே அதற்குத் திறந்த எடுகோள்கள் இல்லை. எனவே அது
தனது இறுதி-!!{formula}~$(!A \land !B) \lif (!B \land !A)$
\emph{இன்} !!a{proof} ஆகும்.
\end{ex}

\end{document}
